Visando facilidade na distribuição e reprodução do ESPEON, optou-se pela criação de um pacote de \textit{software} para navegadores de \textit{internet}, uma extensão Chrome. A escolha se deu pelo suporte ofertado pela Chrome API, que permite --- mediante consentimento --- fazer o monitoramento discente, com acesso à permissões de mídia e conteúdo acessado. A extensão é escrita em JavaScript e tem como componente indispensável o arquivo de manifesto, que contém metadados sobre a extensão, permissões necessárias e \textit{scripts} usados. 

A publicação do pacote de \textit{software} deve ser realizada por meio de uma conta de desenvolvedor. Contudo, embora fosse possível disponibilizar o ESPEON na Chrome Web Store, optou-se pela distribuição local para reduzir a exposição e mitigar potenciais acessos maliciosos ao \textit{software}.

\subsection{Distribuição Local}\label{sub:espeon_distro}
A estratégia de distribuição local consiste na aquisição do código-fonte pelo usuário, seguida do carregamento do pacote descompactado na aba de extensões do navegador, usando o modo desenvolvedor.

Primeiramente, o usuário acessa o código-fonte no repositório GitHub e faz uma cópia local a partir da \textit{branch main}, seja por clone ou \textit{download}. Depois disso, deve-se acessar o navegador na aba de extensões e carregar o pacote, conforme indicado na Figura \ref{fig:load_unpacked}.

\begin{figure}[H]
    \centering
    \begin{minipage}{0.9\textwidth}
        \centering
        \includegraphics[width=.9\textwidth]{espeon/load.unpacked.png}
        \caption{Carregando o Pacote}
        \label{fig:load_unpacked}
    \end{minipage}\hfill
\end{figure}

Carregada a extensão, pode-se habilitá-la no \textit{checkbox} do canto superior direito e depois acessá-la na aba correlata, como ilustrado na Figura \ref{fig:start_extension}.

\begin{figure}[H]
    \centering
    \begin{minipage}{0.9\textwidth}
        \centering
        \includegraphics[width=.9\textwidth]{espeon/iniciar.extensao.png}
        \caption{Iniciando Extensão}
        \label{fig:start_extension}
    \end{minipage}\hfill
\end{figure}

Dessa forma, o ESPEON já está habilitado para uso, cabendo ao aluno fazer seu registro na plataforma para autenticação e posterior ingresso em webconferência (aula).

O processo descrito pode ser sumarizado conforme o diagrama de fluxo disposto na Figura~\ref{fig:flowchart_espeon_intall} a seguir.

\begin{figure}[H]
    \centering
    \begin{minipage}{0.9\textwidth}
        \centering
        \includegraphics[width=.9\textwidth]{draw.io/instalacao.espeon.drawio.png}
        \caption{Diagrama de Fluxo - Instalação ESPEON}
        \label{fig:flowchart_espeon_intall}
    \end{minipage}\hfill
\end{figure}

Trata-se de um processo a ser desempenhado uma única vez; há persistência da extensão na aba correlata do navegador a menos que o usuário opte por removê-la.